\problem{P37}{固定位置的近邻距离}{邻域计数}
在单位线段上独立均匀抽取 $n\ge1$ 个样本点，另固定一个观测位置 $a\in[0,1]$，它不属于随机样本。把样本到 $a$ 的距离从小到大排列，记第 $k$ 个为 $D_{(k)}$，$1\le k\le n$。求一般 $a$ 下 $D_{(k)}$ 的分布，并在 $a=0$ 与 $a=1/2$ 时给出密度、均值和方差。
\solution 对 $r\ge0$，观测位置半径 $r$ 的邻域长度为
\[
\ell_a(r)=\min(1,a+r)-\max(0,a-r).
\]
这里当 $r$ 超过所有可能距离时，长度为 $1$。每个样本独立以概率 $\ell_a(r)$ 落入邻域，故邻域样本数为参数 $n,\ell_a(r)$ 的二项变量。第 $k$ 近距离超过 $r$，恰好等价于邻域中至多 $k-1$ 个样本。因此
\[
\boxed{\Prob(D_{(k)}>r)=\sum_{j=0}^{k-1}\binom nj
\ell_a(r)^j[1-\ell_a(r)]^{n-j},\quad r\ge0.}
\]
当 $a=0$ 时，单个距离本身均匀分布于 $(0,1)$，于是
\[
f_{D_{(k)}}(r)=\frac{n!}{(k-1)!(n-k)!}r^{k-1}(1-r)^{n-k},\quad0<r<1.
\]
这是 $\operatorname{Beta}(k,n-k+1)$ 密度，贝塔积分给出
\[
\boxed{\E D_{(k)}=\frac{k}{n+1},\qquad
\Var D_{(k)}=\frac{k(n-k+1)}{(n+1)^2(n+2)}.}
\]
当 $a=1/2$ 时，单个距离的分布函数在 $[0,1/2]$ 上为 $2r$，故 $2D_{(k)}$ 服从同一贝塔分布。相应密度为
\[
f_{D_{(k)}}(r)=\frac{2n!}{(k-1)!(n-k)!}(2r)^{k-1}(1-2r)^{n-k},\quad0<r<1/2,
\]
均值为端点情形的一半，方差为四分之一。
\insight 把第 $k$ 近距离超过阈值的事件转成固定邻域中的计数事件，可统一处理任意位置和任意次序。端点与中点的区别来自邻域长度的增长速度，而不是样本之间的互相配对。
